\chapter*{TÀI LIỆU THAM KHẢO}
\addcontentsline{toc}{chapter}{TÀI LIỆU THAM KHẢO}

\begin{enumerate}
    \item \textbf{ReactJS Official Documentation}: Nền tảng kiến thức cốt lõi về React, vòng đời Component, và Hooks. Truy cập tại: \url{https://react.dev/}
    \item \textbf{Zustand Documentation}: Hướng dẫn chi tiết về thư viện quản lý trạng thái nhỏ gọn, linh hoạt và cách tích hợp Middleware. Truy cập tại: \url{https://docs.pmnd.rs/zustand/}
    \item \textbf{dnd-kit (Drag \& Drop toolkit)}: Tài liệu kỹ thuật cốt lõi về cơ chế kéo thả, cấu hình Sensors, Modifiers và các thuật toán Collision Detection. Truy cập tại: \url{https://docs.dndkit.com/}
    \item \textbf{React Router v6 Guide}: Cách triển khai định tuyến linh hoạt, xây dựng Protected Routes và Data-fetching. Truy cập tại: \url{https://reactrouter.com/}
    \item \textbf{Tailwind CSS Documentation}: Thư viện utility-first CSS hỗ trợ xây dựng giao diện phức tạp và tối ưu Responsive. Truy cập tại: \url{https://tailwindcss.com/docs}
    \item \textbf{Immer}: Kỹ thuật xử lý cập nhật trạng thái lồng nhau (Deep nested state) bằng cách sử dụng các đối tượng bất biến (Immutable). Truy cập tại: \url{https://immerjs.github.io/immer/}
    \item \textbf{Feature-Sliced Design (FSD)}: Phương pháp luận và nguyên tắc thiết kế cấu trúc hệ thống kiến trúc Frontend dành cho các dự án quy mô lớn. Truy cập tại: \url{https://feature-sliced.design/}
    \item \textbf{ViteJS}: Tìm hiểu về Build tool thế hệ mới dựa trên esbuild và quy trình cấu hình đóng gói mã nguồn. Truy cập tại: \url{https://vitejs.dev/guide/}
    \item \textbf{MDN Web Docs - Web Storage API}: Bách khoa toàn thư từ Mozilla về cơ chế hoạt động của LocalStorage và cấu trúc JSON. Truy cập tại: \url{https://developer.mozilla.org/en-US/docs/Web/API/Web_Storage_API}
    \item \textbf{Vitest \& React Testing Library}: Các tài liệu tham khảo về phương pháp kiểm thử Đơn vị (Unit Test) và kiểm thử Giao diện (Component Test). Truy cập tại: \url{https://vitest.dev/} và \url{https://testing-library.com/}
\end{enumerate}

\appendix

\chapter{Nhật ký công việc (Weekly Log Summary)}

\section*{Tuần 1: Khởi động và Nghiên cứu}
\begin{itemize}
    \item \textbf{Mai Trần Tuấn Kiệt:} Tìm hiểu quy trình làm việc của Trello, khởi tạo dự án với Vite và ReactJS. -- \textit{Hoàn thành}
    \item \textbf{Vũ Lê Đức Anh:} Nghiên cứu các thư viện hỗ trợ: TailwindCSS để làm giao diện và Lucide React cho icon. -- \textit{Hoàn thành}
    \item \textbf{Phạm Nguyễn Nguyên Khôi:} Tìm hiểu về cấu trúc dữ liệu JSON để lưu trữ Board, List và Card. -- \textit{Hoàn thành}
\end{itemize}

\section*{Tuần 2: Kiến trúc ứng dụng và Routing}
\begin{itemize}
    \item \textbf{Mai Trần Tuấn Kiệt:} Thiết lập React Router, tạo các trang trống cho Login, Dashboard và Board Detail. -- \textit{Hoàn thành}
    \item \textbf{Vũ Lê Đức Anh:} Xây dựng Layout chung cho ứng dụng (Navbar, Sidebar cơ bản) và thiết lập cấu trúc Feature-Sliced Design. -- \textit{Hoàn thành}
    \item \textbf{Phạm Nguyễn Nguyên Khôi:} Thiết lập Zustand Store cơ bản để chuẩn bị quản lý Global State. -- \textit{Hoàn thành}
\end{itemize}

\section*{Tuần 3: Authentication và Bảo mật}
\begin{itemize}
    \item \textbf{Mai Trần Tuấn Kiệt:} Cấu hình Zustand Slice để quản lý trạng thái đăng nhập và thông tin phiên làm việc. -- \textit{Hoàn thành}
    \item \textbf{Vũ Lê Đức Anh:} Thiết kế giao diện trang Login trang trọng và trực quan. -- \textit{Hoàn thành}
    \item \textbf{Phạm Nguyễn Nguyên Khôi:} Xử lý logic Protected Routes (Chặn truy cập khi chưa đăng nhập). -- \textit{Hoàn thành}
\end{itemize}

\section*{Tuần 4: Quản lý Board và Dashboard}
\begin{itemize}
    \item \textbf{Mai Trần Tuấn Kiệt:} Xây dựng trang Dashboard hiển thị danh sách các Board hiện có. -- \textit{Hoàn thành}
    \item \textbf{Vũ Lê Đức Anh:} Thêm chức năng tạo mới Board và xóa Board khỏi danh sách. -- \textit{Hoàn thành}
    \item \textbf{Phạm Nguyễn Nguyên Khôi:} Viết Zustand \texttt{boardSlice} để xử lý logic lưu trữ Board vào Global State. -- \textit{Hoàn thành}
\end{itemize}

\section*{Tuần 5: Thành phần tĩnh của Board}
\begin{itemize}
    \item \textbf{Mai Trần Tuấn Kiệt:} Thiết kế giao diện tĩnh cho Column (List) và Card bên trong Board. -- \textit{Hoàn thành}
    \item \textbf{Vũ Lê Đức Anh:} Xây dựng form thêm nhanh List mới và Card mới. -- \textit{Hoàn thành}
    \item \textbf{Phạm Nguyễn Nguyên Khôi:} Xử lý logic hiển thị dữ liệu từ State lên giao diện Board Detail. -- \textit{Hoàn thành}
\end{itemize}

\section*{Tuần 6: Hoàn thiện UI và Cấu trúc Component}
\begin{itemize}
    \item \textbf{Mai Trần Tuấn Kiệt:} Tách các component lớn (\texttt{Dashboard}, \texttt{BoardDetail}) thành các component nhỏ hơn (\texttt{BoardCard}, \texttt{List}, \texttt{TaskCard}) để tối ưu tái sử dụng. -- \textit{Hoàn thành}
    \item \textbf{Vũ Lê Đức Anh:} Xây dựng \texttt{CardModal} để hiển thị chi tiết thẻ (nhãn, ngày hết hạn, mô tả). -- \textit{Hoàn thành}
    \item \textbf{Phạm Nguyễn Nguyên Khôi:} Loại bỏ các thành phần UI thừa, viết tài liệu hướng dẫn nội bộ (\texttt{docs}) để chuẩn bị tích hợp logic kéo thả. -- \textit{Hoàn thành}
\end{itemize}

\section*{Dự kiến tuần 7 - 9:}
\begin{itemize}
    \item \textbf{Tuần 7:} Tích hợp thư viện \texttt{dnd-kit} để xử lý kéo thả Card và List.
    \item \textbf{Tuần 8:} Cài đặt Local Storage để lưu dữ liệu bền vững và tối ưu hóa hiệu năng.
    \item \textbf{Tuần 9:} Kiểm thử hệ thống, sửa lỗi và hoàn thiện báo cáo toàn văn.
\end{itemize}

\chapter{Minh chứng làm việc nhóm (Git History)}
[Chèn ảnh chụp màn hình lịch sử commit và biểu đồ đóng góp của các thành viên tại đây]
